\section{A worked analysis of a recoverable repository workflow}\label{sec:C20:case}
Consider a repository-maintenance agent asked to update a dependency, run tests, and prepare a reviewable change. Its environment has four independently mutable surfaces: the checkout, dependency registry, test infrastructure, and remote review service. The task is deliberately narrower than autonomous deployment. The relevant success predicate requires the intended version change, passing required tests under the declared environment, and a review artifact referring to the verified revision. A passing local command is not sufficient if the remote artifact contains another commit.

Start by assigning an immutable task identifier and recording the initial repository commit. The package resolution step produces a lockfile whose content depends on registry state as well as the requested version. Therefore the transition record must include the resolved artifact identity, not only the human-readable version string. A registry can republish metadata or a moving tag can resolve differently. Replaying the same command later may produce different bytes even when the agent's prompt is unchanged. Reproducibility requires controlling this dependency or recording that it remains uncontrolled.

Next, the agent edits files in a workspace branch. The harness can checkpoint a filesystem snapshot before execution, but the snapshot does not capture every external dependency. A test can consult network services, local caches, or environment variables. The verification record should bind the resulting verdict to the workspace digest, test selection, dependency manifest, and relevant environment identity. If the model changes a source file after tests finish, the prior verdict becomes stale. The runtime can enforce this mechanically by invalidating verdicts whose dependency digests no longer match the current artifact.

This binding can be expressed as an acceptance rule. Let $x$ be the proposed artifact, $d(x)$ its digest, $v$ a verifier configuration, and $r$ a recorded test outcome. The gate accepts only if $r$ was produced for $(d(x),v)$ and the required checks passed. The rule establishes correspondence between artifact and evidence; it does not establish completeness of the tests. An independent specification defect remains possible. The separation is useful because it prevents a simple engineering mismatch from being obscured by philosophical uncertainty about correctness.

Now introduce a crash after the remote review service creates the artifact but before the harness stores its identifier. Replaying the create request can produce duplicates. If the service supports an idempotency key, recovery uses the original key. If it does not, the agent can query for an artifact marked with the task identifier and expected commit digest. That query may return zero, one, or several matches. Zero results are not conclusive under eventual consistency. Several results require reconciliation rather than arbitrary selection. A bounded waiting policy and an unresolved status are more accurate than asserting that the operation failed.

The cost of naive retries can be analyzed without inventing a measurement. Conditional on one dispatch, let $A$ be the event that the request takes effect but its acknowledgment is lost, and let $R$ mean recovery retries without reconciliation. For $q=\Pr(A)>0$ and $r=\Pr(R\mid A)$, the chain rule gives $\Pr(A\cap R)=qr$. No independence assumption is needed with this conditional definition of $r$. If $q=0$, the joint event has probability zero without defining $r$. Under the one-retry model this is a duplicate opportunity. The actual duplicate probability also depends on remote deduplication and whether the operation is intrinsically idempotent. This expression is not an estimate of any product. It identifies two distinct intervention targets: reduce acknowledgment ambiguity or eliminate blind retries. Improving model reasoning alone need not reduce either one.

Suppose authorization changes while the task is interrupted. The original grant allowed preparation of a review artifact, but the operator later revokes remote write access. Restoring the old checkpoint must not restore the old grant as if it were still current. Recovery loads execution progress from the checkpoint and authorization from the authoritative policy service. The resumed model may still see a historical permission statement in $H_t$; the tool boundary must reject a newly attempted write under the current epoch. This example shows why replay fidelity and present authorization can legitimately diverge.

Deletion introduces another dependency. If a user requests removal of sensitive task material, deleting the current context does not remove journal payloads, filesystem checkpoints, test logs, remote artifacts, or indexing copies. A retention service needs a dependency inventory and a distinction between records required for audit and content permitted for reuse. Where legal or operational retention prevents immediate physical deletion, the system should record the constraint and prevent subsequent model retrieval. It should not describe a tombstone as completed physical erasure. Chapter~\ref{ch:C29} examines the authority and governance implications.

A controlled evaluation of this workflow uses deterministic repository fixtures and instrumented remote-service simulators. Each fixture specifies the allowed edits, required checks, and permitted remote effects. Failures are injected at journal, network, verifier, and artifact-publication boundaries. The primary unit is a task instance with its complete failure schedule, not an individual model call. Paired schedules allow comparison of recovery strategies without confounding one strategy with easier failures. Model sampling remains a separate repeated-run factor, and uncertainty should preserve clustering by fixture.

Report artifact correctness, stale-verdict acceptance, duplicate creations, authorization violations, unresolved outcomes, and complete execution cost. A runtime can achieve fewer duplicates by refusing all remote work, so benign completion must be reported alongside containment. It can achieve low mean latency by abandoning difficult recoveries, so abandoned and timed-out runs remain visible. These paired metrics force the experiment to evaluate the desired tradeoff rather than reward a degenerate implementation. The worked example yields no universal winner, but it supplies a precise testable distinction between durable progress and merely persistent conversation.
